\documentclass[11pt]{article}
\usepackage[utf8]{inputenc}
\usepackage[T1]{fontenc}
\usepackage[margin=2.5cm]{geometry}
\usepackage{amsmath,amssymb,amsthm,mathtools}
\usepackage{enumitem}
\usepackage{xcolor}
\usepackage{hyperref}
\hypersetup{colorlinks=true,linkcolor=blue!60!black,urlcolor=blue!60!black}

\theoremstyle{plain}
\newtheorem{theorem}{Theorem}[section]
\newtheorem{proposition}[theorem]{Proposition}
\newtheorem{lemma}[theorem]{Lemma}
\newtheorem{corollary}[theorem]{Corollary}
\theoremstyle{definition}
\newtheorem{assumption}{Assumption}
\theoremstyle{remark}
\newtheorem{remark}[theorem]{Remark}

\newcommand{\R}{\mathbb{R}}
\newcommand{\N}{\mathbb{N}}
\newcommand{\E}{\mathbb{E}}
\newcommand{\KL}{\mathrm{KL}}
\DeclareMathOperator*{\argmin}{arg\,min}
\DeclareMathOperator{\dimm}{dim}

\title{\textbf{Research Plan --- Minimax Estimation on Orbit Spaces:\\
the Exact Rate Gain from Invariance via Orbital MDL\\
and Concentration in the Multiplicative Twin Space}}
\author{J.~Nemb\'e \\ \small (plan / extended abstract --- not for submission as is)}
\date{\today}

\begin{document}
\maketitle

\begin{abstract}
This is a \emph{plan}, not a finished article: it fixes the target theorems and
the proof architecture, reusing the building blocks already validated and
corrected (orbital Glivenko--Cantelli/Donsker; Barron--Cover/Hellinger oracle
inequality; Fano lower bound). The thesis is that orbit-based inference does
carry genuine added value, but \emph{not} where it was originally placed. A
global gauge conjugation is an isomorphism and cannot improve a convergence
theorem; the value appears only when the orbit structure genuinely reduces the
\emph{effective dimension} of the problem. We make this precise: for a
$G$-invariant target on $E$, the minimax risk is governed by
$d_{\mathrm{eff}}=\dimm(E/G)=\dimm E-\dimm(\text{orbit})$, and is achieved
adaptively by an \emph{orbital minimum-description-length} estimator. The three
tools of the program are slotted in their correct roles: (1) data complexity
$\to$ \emph{metric entropy of the quotient}; (2) candidate complexity $\to$
\emph{orbital approximation / description length}; (3) concentration in twin
spaces $\to$ the \emph{multiplicative} (logarithmic-twin) space of Hellinger
affinities, where the Barron--Cover partition-function bound is a genuine
large-deviation estimate. The \emph{distinctive contribution} we target is the
combination of the rate gain with a \emph{multiplicative/heavy-tailed} data
regime (positive intensities, compositional data): there the affinity coordinate
is canonical and scale-invariant, and the orbital MDL estimator attains the same
gained rate where identity-coordinate concentration fails --- a regime outside
the existing invariance-gain literature. A frank novelty assessment versus the
recent ``sample-complexity gain from invariance'' literature closes the plan.
\end{abstract}

\tableofcontents

\section{The conceptual correction (why the program is sound but was misplaced)}

A twin space $(E,\odot_\phi)$ is the image of a classical structure under a
bijection $\phi$ (the Transfer Theorem). Hence any concentration/large-deviation
statement ``in the twin space'' is the exact push-forward of a classical one:
conjugation by a bijection is an isomorphism and creates no information, tightens
no bound, and cannot lower the effective sample size below the number of
independent observations. This single fact explains every error corrected in the
source notes (the spurious $\sqrt{m}$ sample inflation; the unconditional KL
oracle inequality; the $\sup$ orbital metric).

\medskip
\noindent\textbf{Where the value actually is.} Two places, both used below:
\begin{enumerate}[label=(\roman*)]
\item \emph{Dimension reduction.} When $G$ acts with \emph{positive-dimensional}
orbits, the invariant problem genuinely lives on the lower-dimensional quotient
$E/G$. The minimax rate then improves from $n^{-2s/(2s+D)}$ to
$n^{-2s/(2s+d_{\mathrm{eff}})}$. For a \emph{finite} group the orbits are
$0$-dimensional, $d_{\mathrm{eff}}=D$, and the gain is only a constant factor ---
exactly consistent with the corrected orbit-averaging theorem (variance reduced
by at most $|G|$, never the rate).
\item \emph{The right twin space for concentration.} The Hellinger affinity
$\rho(P,Q)=\int\sqrt{dP\,dQ}$ is \emph{multiplicative} over products,
$\rho(P^{\otimes n},Q^{\otimes n})=\rho(P,Q)^n$. Thus $-\log\rho$ is additive:
the natural ``twin sum'' of evidence is the multiplication of likelihood ratios,
i.e.\ addition in the logarithmic-twin coordinate. The Barron--Cover/Kraft
partition-function bound is precisely a Cram\'er-type large-deviation estimate in
\emph{this} twin space --- not in the data space. This is the legitimate form of
the third pillar.
\end{enumerate}

\section{Setting and standing assumptions}

\begin{assumption}[Geometry and group]\label{ass:geo}
$(E,\rho_E)$ is a compact metric space (typically a compact Riemannian manifold
of dimension $D$). A compact group $G$ acts continuously on $E$ by isometries,
with closed orbits; generic orbits have dimension $r:=\dimm(G\cdot x)$. The
quotient $E/G$ is a compact metric space for the \emph{orbital (infimum) metric}
\[
d_G([x],[y]) := \inf_{g\in G}\rho_E(x,\,g\cdot y),
\]
(the corrected metric --- the supremum form is not well defined on the quotient),
with covering dimension $d_{\mathrm{eff}}:=\dimm(E/G)=D-r$ (for a smooth proper
action).
\end{assumption}

\begin{assumption}[Reference measure and invariant class]\label{ass:class}
Fix a $G$-invariant base measure $\mu$ on $E$ with push-forward $\mu_G:=\pi_*\mu$
on $E/G$, $\pi:E\to E/G$. The target is \emph{known to be $G$-invariant},
$f^\star=\tilde f^\star\circ\pi$, and $\tilde f^\star$ lies in the H\"older (or
Sobolev) ball $\Sigma(s,L)$ of order $s>0$, radius $L$, on $(E/G,d_G,\mu_G)$. We
write $\Sigma_{\mathrm{inv}}(s,L)$ for the pulled-back invariant class on $E$.
\end{assumption}

\begin{remark}
Assumption~\ref{ass:class} (the target is invariant) is the \emph{structural
prior} that makes the gain real; without it, orbit projection estimates only the
invariant part $\pi_*P^\star$ and there is nothing to gain.
\end{remark}

\section{Tool 1 --- data complexity as metric entropy of the quotient}

\begin{lemma}[Orbital covering numbers]\label{lem:cov}
Under Assumptions~\ref{ass:geo}--\ref{ass:class}, the invariant class obeys
\[
\log N\big(\varepsilon,\ \Sigma_{\mathrm{inv}}(s,L),\ \|\cdot\|_\infty\big)
\ \asymp\ \Big(\tfrac{L}{\varepsilon}\Big)^{d_{\mathrm{eff}}/s},
\]
governed by $d_{\mathrm{eff}}=\dimm(E/G)$, not by $D$.
\end{lemma}

\noindent\emph{Proof sketch.} A $G$-invariant function on $E$ is a function on
$E/G$; cover $(E/G,d_G)$ by $\asymp(1/\varepsilon)^{d_{\mathrm{eff}}}$ balls and
apply the standard entropy bound for H\"older balls on a
$d_{\mathrm{eff}}$-dimensional metric space. This is the corrected covering-number
lemma of the orbital-estimation note, restricted to the quotient. \hfill$\square$

\section{Tool 2 --- candidate complexity by orbital approximation (description length)}

\begin{lemma}[Orbital sieve and twin Jackson bound]\label{lem:sieve}
There is a countable family of $G$-invariant models $\{\mathcal M_m\}_{m\in\mathcal I}$
(invariant histograms/wavelets at resolution $2^k$ pulled back from $E/G$;
$m=(k,\theta)$) with
\[
\dimm\mathcal M_k\ \asymp\ 2^{k\,d_{\mathrm{eff}}},\qquad
\inf_{q\in\mathcal M_k}\KL(P^\star\|Q)\ \le\ C\,L^2\,2^{-2ks}
\quad(f^\star\in\Sigma(s,L)),
\]
and description lengths $L(m)$ satisfying Kraft $\sum_m 2^{-L(m)}\le1$ with
$L(k,\theta)\asymp \dimm\mathcal M_k\cdot\log(\cdot)$.
\end{lemma}

\noindent\emph{Proof sketch.} Pull a standard linear sieve on $E/G$ back through
$\pi$; invariance is automatic. The approximation (Jackson) bound is the
classical one on the $d_{\mathrm{eff}}$-dimensional quotient. Coding the
resolution $k$ and the cell counts gives a prefix code; replace $L$ by $2L$ to
secure the half-exponent Kraft condition $\sum 2^{-L(m)/2}\le1$ used in
Tool~3. This is the orbital realization of ``choosing the candidate function from
the orbit structure.'' \hfill$\square$

\section{Tool 3 --- concentration in the multiplicative twin space}

\begin{proposition}[Affinity is the twin sum of evidence]\label{prop:affinity}
For the Hellinger affinity $\rho(P,Q)=\int\sqrt{dP\,dQ}\,d\mu$,
$\rho(P^{\otimes n},Q^{\otimes n})=\rho(P,Q)^n$; equivalently
$-\log\rho(P^{\otimes n},Q^{\otimes n})=n\,(-\log\rho(P,Q))$. Concentration of
log-likelihood ratios is therefore additive in the logarithmic-twin coordinate,
and the Kraft partition-function bound
\[
\E_{P^\star}\!\Big[2^{-L(\widehat m)/2}\textstyle\prod_{i=1}^n
\sqrt{\widehat f(X_i)/f^\star(X_i)}\Big]\ \le\ \sum_m 2^{-L(m)/2}\,\rho(f^\star,f_m)^n\ \le\ 1
\]
is a Cram\'er-type large-deviation estimate in the multiplicative twin space
$(\R_{>0},\times)$.
\end{proposition}

This is the corrected, well-posed third pillar: the relevant non-additive
structure is the multiplicative space of affinities/likelihood ratios, where the
``twin sum'' genuinely is multiplication.

\begin{remark}[Why the affinity, and not a free coordinate choice]\label{rem:nofreelunch}
The companion analysis (\texttt{concentration\_diffeomorphism\_corrected})
establishes two facts that make Proposition~\ref{prop:affinity} the \emph{only}
legitimate form of the third pillar. First, \emph{no free lunch}: for a fixed
functional the bounded-difference/sub-Gaussian constants are
\emph{coordinate-intrinsic}, so a gauge change cannot tighten a concentration
bound (and the unnormalised optimum $\inf_\phi F[\mu,\phi]=0$ is degenerate).
Hence concentration cannot, by itself, improve a convergence rate --- which is
exactly why the rate gain here comes from dimension reduction (Tools~1--2), not
from Tool~3. Second, the affinity $\rho$ is the canonical \emph{scale-invariant}
concentration object: $-\log\rho$ is additive and unit-free, curing the
scale-dependence that invalidated coordinate ``optimisation''. Tool~3 therefore
serves to make the oracle inequality self-contained, not to create rate.
\end{remark}

\section{Main results (targets)}

\begin{theorem}[Upper bound: adaptive orbital MDL]\label{thm:upper}
Let $\widehat f=f_{\widehat m}$ be the orbital MDL estimator
\[
\widehat m\in\argmin_{m\in\mathcal I}\Big\{-\textstyle\sum_{i=1}^n\log f_m(X_i)+L(m)\log2\Big\}.
\]
Then, for every $f^\star\in\Sigma_{\mathrm{inv}}(s,L)$ and every $s\in(0,s_{\max}]$,
\[
\E\big[h^2(P^\star,\widehat f)\big]
\ \le\ C\,\inf_{k}\Big\{L^2\,2^{-2ks}+\tfrac{2^{k\,d_{\mathrm{eff}}}\log n}{n}\Big\}
\ \le\ C'\,\Big(\tfrac{\log n}{n}\Big)^{2s/(2s+d_{\mathrm{eff}})},
\]
\emph{adaptively} (the estimator does not use $s$); the $\log n$ is the two-stage
description-length / adaptation cost, and is removed (exact rate, adaptively) by
the projection estimators of \texttt{orbital\_upper\_bound\_proof}~\S6.
\end{theorem}

\noindent\emph{Proof sketch.} The corrected Barron--Cover argument
(Proposition~\ref{prop:affinity}) yields the Hellinger oracle inequality
$\E[h^2(P^\star,\widehat f)]\le C\inf_m\{\inf_{q\in\mathcal M_m}\KL(P^\star\|q)+\log2\,L(m)/n\}$.
Insert the bias (Lemma~\ref{lem:sieve}) and $L(m)/n\asymp 2^{k d_{\mathrm{eff}}}/n$,
then optimise $k$: $2^{k^\star}\asymp n^{1/(2s+d_{\mathrm{eff}})}$. Adaptivity is
automatic because the penalised criterion selects $k$. \hfill$\square$

\begin{theorem}[Lower bound: Fano on the quotient]\label{thm:lower}
\[
\inf_{\widehat f}\ \sup_{f^\star\in\Sigma_{\mathrm{inv}}(s,L)}
\E\big[h^2(P^\star,\widehat f)\big]\ \ge\ c\,n^{-2s/(2s+d_{\mathrm{eff}})}.
\]
\end{theorem}

\noindent\emph{Proof sketch.} Build $M\asymp\exp(2^{k d_{\mathrm{eff}}})$ invariant
perturbations $f_\omega=f_0\big(1+\delta\sum_{j}\omega_j\,\psi_{k,j}\big)$ with
$\omega\in\{\pm1\}^{N}$, $N\asymp2^{k d_{\mathrm{eff}}}$, where the bumps
$\psi_{k,j}$ are supported on disjoint balls of the \emph{quotient}
$(E/G,d_G)$ pulled back to $E$ (hence automatically $G$-invariant, by the
\emph{infimum} metric of Assumption~\ref{ass:geo}). A Varshamov--Gilbert subset
gives pairwise Hellinger separation $\beta\asymp\delta^2 2^{-k d_{\mathrm{eff}}}\cdot N\asymp\delta^2$
and pairwise KL $\le\alpha$ with $n\alpha\asymp n\delta^2$. Using the
\emph{corrected} Fano bound (the $n\alpha$ factor, not $\alpha$), balance
$\delta^2\asymp 2^{-2ks}$ and $2^{k}\asymp n^{1/(2s+d_{\mathrm{eff}})}$. \hfill$\square$

\begin{theorem}[Exact rate gain from invariance]\label{thm:main}
Under Assumptions~\ref{ass:geo}--\ref{ass:class},
\[
\boxed{\ \mathcal R_n^\star\big(\Sigma_{\mathrm{inv}}(s,L)\big)\ \asymp\ n^{-2s/(2s+d_{\mathrm{eff}})},
\qquad d_{\mathrm{eff}}=\dimm(E/G)=D-r.\ }
\]
The invariant rate strictly improves on the non-invariant rate
$n^{-2s/(2s+D)}$ \emph{iff} the orbits are positive-dimensional ($r>0$); for
finite $G$ ($r=0$) the gain is only a constant factor.
\end{theorem}

\begin{remark}[Exact rate, adaptively: the $\log$ is removed]
$\mathcal R_n^\star$ denotes the minimax rate, and the box is now an honest
two-sided match. The lower bound (Theorem~\ref{thm:lower}, proved in
\texttt{orbital\_lower\_bound\_proof}) is $n^{-2s/(2s+d_{\mathrm{eff}})}$ exactly.
The \emph{adaptive} orbital MDL of Theorem~\ref{thm:upper} attains it up to a
$\log n$ two-stage/coding cost; but this $\log$ is \emph{not} intrinsic and is
removed in \texttt{orbital\_upper\_bound\_proof} (\S6): a single-resolution
\emph{projection} estimator attains $n^{-2s/(2s+d_{\mathrm{eff}})}$ exactly for
known $s$, and a \emph{penalised projection} estimator (Birg\'e--Massart) attains
it \emph{adaptively}. The reason is structural --- a projection / penalised
estimator pays the model's metric complexity $D_k/n\asymp2^{kd_{\mathrm{eff}}}/n$,
whereas the MDL code additionally pays $\log(1/\delta)\asymp\log n$ bits to
quantise each coefficient; the geometric growth $D_k\asymp2^{kd_{\mathrm{eff}}}$
keeps the selection weights summable with linear-in-$D_k$ exponents, so no $\log$
re-enters. Hence the box holds at the \emph{exact} rate, both for known $s$ and
adaptively.
\end{remark}

\section{Examples calibrating the gain (and the honest negative cases)}

\begin{enumerate}[label=(\alph*)]
\item \textbf{Radial densities.} $E=\R^D$ (or $S^D$), $G=SO(D)$:
$d_{\mathrm{eff}}=1$, rate $n^{-2s/(2s+1)}$ vs.\ $n^{-2s/(2s+D)}$ --- maximal gain.
\item \textbf{Partial translation invariance.} $E=\mathbb T^D$, $G=\mathbb T^k$:
$d_{\mathrm{eff}}=D-k$.
\item \textbf{Finite symmetry (negative control).} Compositional data on the
simplex with $G=S_d$, or $m$-fold rotational symmetry $G=\mathbb Z_m$: $r=0$,
\emph{no rate change}, only a constant ($1/|G|$) improvement --- exactly the
corrected orbit-averaging statement. This example is kept deliberately, to show
the theory predicts \emph{when there is no gain}.
\item \textbf{Deformed-operation variant.} On positive/compositional data the
natural ``addition'' is a twin operation $\odot$ (log-ratio, Einstein addition);
$\Sigma(s,L)$ is then smoothness in the twin coordinate. The same theorems hold
with the gauge built into $\pi$, illustrating the non-additive reach.
\end{enumerate}

\section{The multiplicative / heavy-tailed extension (the distinctive contribution)}
\label{sec:multiplicative}

This is the angle that sets the paper apart from the existing invariance-gain
literature, which works in $L^2$/Sobolev/RKHS over Euclidean or bounded domains.
Here the data are \emph{positive} or \emph{compositional}, the natural group
acts through a twin (non-additive) operation, and the identity coordinate admits
no usable concentration; yet the gained rate persists.

\begin{assumption}[Multiplicative regime]\label{ass:mult}
$E\subseteq(0,\infty)^D$ (or the simplex $\Delta^{D-1}$), with the gauge
$\phi=\log$ (resp.\ the centred log-ratio) so that the base operation is the twin
sum $x\odot y=\phi^{-1}(\phi(x)+\phi(y))$. The group $G$ acts by isometries of the
twin coordinate $\phi(E)$, with orbital metric and quotient as in
Assumption~\ref{ass:geo} computed in the $\phi$-chart; $d_{\mathrm{eff}}=\dimm(\phi(E)/G)$.
Smoothness $\Sigma(s,L)$ is measured in the twin coordinate.
\end{assumption}

\begin{lemma}[Matched-coordinate concentration; from the corrected concentration note]\label{lem:matched}
Let the per-coordinate (log-)ranges be $\log(b_i/a_i)$. Then for the
multiplicative aggregate the deviation of $\log$ concentrates at the
\emph{dimensionless} scale $\sum_i\log^2(b_i/a_i)$, and for heavy-tailed
positive data with sub-Gaussian $\log X$ the affinity bound of
Proposition~\ref{prop:affinity} holds in the $\phi$-chart even when no
identity-coordinate (Hoeffding/Bernstein) bound exists. The improvement over the
identity chart is \emph{qualitative} (finite vs.\ infinite), not a numerical
factor on a common functional.
\end{lemma}

\noindent\emph{Proof sketch.} Apply Theorems ``product''/``heavy'' of
\texttt{concentration\_diffeomorphism\_corrected} in the $\phi$-chart; the
affinity is computed with respect to the twin-transported reference measure, and
Proposition~\ref{prop:affinity} is unchanged because $\rho$ is invariant under
the change of variables. \hfill$\square$

\begin{theorem}[Orbital MDL on multiplicative data]\label{thm:mult}
Under Assumptions~\ref{ass:geo},\ref{ass:class},\ref{ass:mult}, the orbital MDL
estimator built on the $\phi$-chart sieve satisfies, for
$f^\star\in\Sigma_{\mathrm{inv}}(s,L)$,
\[
\E\big[h^2(P^\star,\widehat f)\big]\ \le\ C'\,\Big(\tfrac{\log n}{n}\Big)^{2s/(2s+d_{\mathrm{eff}})},
\]
adaptively, with $d_{\mathrm{eff}}=\dimm(\phi(E)/G)$. The rate is identical to the
additive case (Theorem~\ref{thm:upper}); the contribution is that it holds on
positive/compositional, possibly heavy-tailed, data where identity-coordinate
methods give no bound --- because the affinity oracle step is tail-free
(\texttt{orbital\_upper\_bound\_proof}, Lemma on the partition bound).
\end{theorem}

\noindent\emph{Proof sketch.} Run the proof of Theorem~\ref{thm:upper} verbatim
in the $\phi$-chart: Lemma~\ref{lem:cov} (quotient entropy) and
Lemma~\ref{lem:sieve} (twin Jackson) are stated metrically and transport through
$\phi$; the oracle inequality uses the scale-invariant affinity bound
(Proposition~\ref{prop:affinity}, Lemma~\ref{lem:matched}); optimise $k$. The
lower bound (Theorem~\ref{thm:lower}) transports identically since $h^2$ is
invariant under the change of variables. \hfill$\square$

\begin{remark}[What is genuinely new here]
The rate gain itself (Theorem~\ref{thm:main}) overlaps known results; what is new
is (a) the \emph{adaptive MDL / Hellinger} route, and (b)
Theorem~\ref{thm:mult}: the gain on \emph{multiplicative / heavy-tailed} data via
the canonical affinity coordinate, a regime the additive invariance-gain
literature does not cover. Concrete instances: positive Poisson/point-process
intensities, compositional (simplex) data, and multiplicative-noise diffusions
--- the very models of the corrected MDL paper.
\end{remark}

\section{Paper outline (once the plan is executed)}
\begin{enumerate}[label=\arabic*.]
\item Introduction: invariance as effective-dimension reduction; the
multiplicative-twin view of evidence.
\item Geometry of the quotient and the orbital metric (Assumption~\ref{ass:geo}).
\item Tool 1: quotient entropy (Lemma~\ref{lem:cov}).
\item Tool 2: orbital sieve, twin Jackson, description length (Lemma~\ref{lem:sieve}).
\item Tool 3: affinity concentration (Proposition~\ref{prop:affinity}).
\item Upper bound, adaptivity (Theorem~\ref{thm:upper}).
\item Lower bound (Theorem~\ref{thm:lower}); exact rate (Theorem~\ref{thm:main}).
\item Examples; finite-group negative control.
\item \textbf{Multiplicative/heavy-tailed extension} (Section~\ref{sec:multiplicative}):
the distinctive contribution --- matched affinity coordinate, orbital MDL on
positive/compositional data (Theorem~\ref{thm:mult}).
\item Discussion: deformed operations; relation to invariant kernel methods.
\end{enumerate}

\section{Frank novelty assessment (read before writing)}
\begin{itemize}[leftmargin=1.4em]
\item \textbf{What is \emph{not} new.} The bare statement ``invariance lowers the
effective dimension to $\dimm(E/G)$'' is, in the $L^2$/Sobolev and kernel-ridge
settings, essentially established (Tahmasebi \& Jegelka, \emph{The exact sample
complexity gain from invariances}, NeurIPS 2023; Bietti, Venturi, Bruna; Elesedy
\& Zaidi). Theorem~\ref{thm:main} in those settings would be a reformulation and
should not be sold as new.
\item \textbf{What can be genuinely new.} (i) The \emph{information-theoretic /
MDL} route: an \emph{adaptive} orbital description-length estimator achieving the
gain in \emph{Hellinger} risk, with the rate emerging from the Kraft
partition-function bound rather than from RKHS spectra. (ii) The extension to
\emph{deformed (twin) operations}, where the invariance is with respect to a
non-additive group of operations and smoothness is measured in the twin
coordinate --- a regime outside the cited works. (iii) The unifying
\emph{multiplicative-twin-space} formulation of concentration
(Proposition~\ref{prop:affinity}) tying the program's third pillar to a standard
large-deviation principle.
\item \textbf{Recommendation.} Build the paper around (i)+(ii) with
Theorem~\ref{thm:mult} (Section~\ref{sec:multiplicative}) as the headline:
\emph{adaptive orbital MDL for invariant density estimation on multiplicative /
heavy-tailed data}. Theorem~\ref{thm:main} calibrates the gain (and is positioned
explicitly as a re-derivation in the cited additive setting), while the
multiplicative regime is the part with clear, uncontested added value. The
corrected concentration note supplies the self-contained, scale-invariant Tool~3
(Remark~\ref{rem:nofreelunch}).
\end{itemize}

\section*{Remaining work / risks}
\begin{itemize}[leftmargin=1.4em]
\item Make Lemma~\ref{lem:sieve} fully rigorous on a curved quotient (orbifold
points where the action is not free: $E/G$ may be a stratified space, and
$d_{\mathrm{eff}}$ is the dimension of the top stratum).
\item Confirm the Hellinger oracle constant and the half-Kraft normalisation
(Tool~3) for the orbital sieve.
\item Decide the deformed-operation scope (which gauges $\phi$ keep
$\Sigma(s,L)$ a genuine smoothness class).
\item Literature due diligence on the exact-gain papers to pin the precise delta.
\item Multiplicative regime (Section~\ref{sec:multiplicative}): integrability of
the twin-transported reference measure, and conditions on the tails of $\log X$
ensuring the affinity bound and the twin Jackson estimate both hold.
\end{itemize}

\end{document}
